\documentclass[10pt,a4paper]{amsart}
\usepackage[margin=19mm]{geometry}
\usepackage{amsmath,amssymb,mathtools}
\usepackage[scr=rsfs,cal=euler]{mathalfa}
\usepackage{xfrac}
\usepackage[unicode,hidelinks,bookmarks=false]{hyperref}
\newcommand{\ie}{i.e.\ }
\newcommand{\cf}{cf.\ }
\newcommand{\resp}{respectively\ }
\newcommand{\Subsection}{\subsection}
\providecommand{\nobreakdash}{\nobreak\hbox{-}}
\setlength{\emergencystretch}{2em}
\multlinegap0pt
\usepackage{bm}
\usepackage{bbm}
\usepackage{dsfont}
\usepackage{xspace}
\usepackage{cancel}
\usepackage{mathtools}
\usepackage{amsthm}
\makeatletter
\@ifundefined{aassumption}{\newtheorem{aassumption}{Assumption}}{}
\@ifundefined{proposition}{\newtheorem{proposition}{Proposition}}{}
\@ifundefined{theorem}{\newtheorem{theorem}{Theorem}}{}
\makeatother
\title[An inexact $q$-order regularized proximal Newton method for ]{An inexact $q$-order regularized proximal Newton method for nonconvex composite optimization}
\author{Ruyu Liu, Shaohua Pan, Yitian Qian}
\date{Source: arXiv:2311.06871v4 (CC BY 4.0)}
\begin{document}
\maketitle

\noindent\emph{Source and licence.} An inexact $q$-order regularized proximal Newton method for nonconvex composite optimization. Version arXiv:2311.06871v4 (CC BY 4.0), distributed under the Creative Commons Attribution 4.0 International licence, as recorded in the arXiv licence metadata for this exact version and shown on the abstract page https://arxiv.org/abs/2311.06871v4. Full rights notice: licenses/arxiv-2311.06871-CC-BY-4.0.txt.\par\smallskip

\noindent\textit{Editorial scope.} Counted unit: the Theorem whose source label is \texttt{Theorem-Lconverge} in the article (source0.tex lines 868--877, source label \texttt{Theorem-Lconverge}), delivered together with the article's own complete original proof of it (source0.tex lines 878--964, one \texttt{proof} environment of 87 lines). No mathematics is written, repaired or re-derived: statement and proof are the article's own text line for line. Required context retained uncounted: ass1 (lines 138--144); subprobk (lines 321--323); new-iter (lines 365--370); yk-vk (lines 409--411); ass2 (lines 511--516); prop2-xk (lines 520--531); grad-Lip (lines 589--591); ineq-quad (lines 622--624); prop-dkbound1 (lines 815--821); ykxk-bound (lines 818--820); ykxk-relation (lines 824--826). Every definition, display and label of those lines is retained unchanged. Omitted from this package and not redistributed: 1--137, 145--320, 324--364, 371--408, 412--510, 517--519, 532--588, 592--621, 625--814, 821--823, 827--867, 965--1365. Nothing inside a retained excerpt is omitted or altered: every retained body is delivered exactly as the article's lines are written, so no omission receipt of any kind accompanies this package. Citations: the retained excerpts carry no citation command at all, so no bibliography entry of the article is transcribed and no reference is redistributed. Reference adaptation: none was needed and none was made; every reference inside the delivered unit resolves inside it, so no unresolved reference or citation is produced. Printed slips kept literal: no printed word, symbol, hypothesis or number of the counted statement or of its proof was altered. Layout and class: the article's own class is \texttt{siamart220329} with packages including graphics, lipsum, amsfonts, graphicx, epstopdf, algorithmic, subfigure, multirow, epsfig, booktabs, cases, amsbsy, amsmath, latexsym; the wrapper is the lane's pinned \texttt{amsart} editorial wrapper at 10pt on A4, which restates the theorem-like environments the retained text needs and adds the article's own macro definitions that the retained text needs (none), with extra wrapper packages (bm, bbm, dsfont, xspace, cancel, mathtools). The delivered numbering is the unit's own. Assets: no retained span contains a figure environment, a graphics-inclusion command, a PDF-page inclusion, a table or a TikZ picture; no image file is copied into this package, the package declares no asset dependency and no image file is redistributed with it. Licence: version arXiv:2311.06871v4 of this article is distributed under the Creative Commons Attribution 4.0 International licence, as recorded in the arXiv licence metadata for this exact version and shown on the abstract page https://arxiv.org/abs/2311.06871v4. A full rights notice is delivered at licenses/arxiv-2311.06871-CC-BY-4.0.txt. Characters: every retained excerpt is pure ASCII, so the delivered engine, XeLaTeX with its default Latin Modern text font, reports no missing character.
\begin{aassumption}\label{ass1}
\begin{description}
 \item[(i)] $f$ is a proper lower semicontinuous (lsc) function that is twice continuously differentiable on an open set $\mathcal{O}\supset{\rm cl}\,({\rm dom}\,g)$;
		
 \item[(ii)] $g$ is a proper lsc convex function with ${\rm dom}\,g\ne\emptyset$, and $F$ is lower bounded.
\end{description} 
\end{aassumption}

\begin{equation}\label{subprobk}
 \min_{x\in\mathbb{X}}\Theta_k(x):=f_k(x)+g(x)\ \ {\rm with}\ f_k(x)\!:=\vartheta_k(x)+({L_k}/q)\|x\!-\!x^k\|^q
\end{equation}

		\begin{equation}\label{new-iter}
			x^{k+1}\!:=\left\{\begin{array}{cl}
				y^k &{\rm if}\ F(y^k)<F(y^k\!-\!v^k),\\
				y^k\!-\!v^k &{\rm  otherwise}.
			\end{array}\right.
		\end{equation}

 \begin{equation}\label{yk-vk}
  y^k\!-\!v^k={\rm prox}_{g}(y^k\!-\!\nabla\!f_{k,j_k}(y^k))\ \ {\rm and}\  \ y^k\!-\!\nabla\!f_{k,j_k}(y^k)\in\partial g(y^k\!-\!v^k)\quad\forall k\in\mathbb{N}.
 \end{equation}

\begin{aassumption}\label{ass2}
{\bf(i)} The sequence $\{x^k\}_{k\in\mathbb{N}}$ is bounded; 
	
 \noindent
 {\bf(ii)} $\nabla^2\!f$ is strictly continuous on an open convex set $\mathcal{N}$ with $\mathcal{O}\supset\mathcal{N}\!\supset\omega(x^0)$.
\end{aassumption}

\begin{proposition}\label{prop2-xk}  
 Under Assumptions \ref{ass1}-\ref{ass2}, the following statements hold.  
 \begin{itemize}	
 \item[(i)] There exists a constant $\widehat{L}>0$ such that $L_k\le\widehat{L}$ for all $k\in\mathbb{N}$.
		
 \item[(ii)] $\lim_{k\to\infty}v^k=\lim_{k\to\infty}R_k(y^k)=0$ and $\lim_{k\to\infty}r(x^k)=0$.
		
 \item[(iii)] $\omega(x^0)$ is a nonempty connected compact set, and $\omega(x^0)\subset\mathcal{S}^*$. 
 
 \item[(iv)] For each $\overline{x}\in\omega(x^0)$, there exists a subsequence $\{x^{k_{\nu}}\}_{\nu\in\mathbb{N}}$ such that $x^{k_{\nu}}\to\overline{x}$ and $F(x^{k_{\nu}})\to F(\overline{x})$ as $\nu\to\infty$, and hence $F(\overline{x})=\lim_{k\to\infty}F(x^k)\!:=\overline{F}$.
\end{itemize}
\end{proposition}

 \begin{equation}\label{grad-Lip}
  \|\nabla\!f(y)-\nabla\!f(x)\|\le L_{g}\|y-x\|\quad\forall x,y\in\Gamma.
 \end{equation}

 \begin{equation}\label{ineq-quad}
 \|\nabla\!f(y^k)-\nabla\!f(x^k)-\nabla^2\!f(x^k)(y^k\!-\!x^k)\|\le({L_H}/{2})\|y^k\!-\!x^k\|^{2}.
 \end{equation}  

\begin{proposition}\label{prop-dkbound1}
 Fix any $q\in\!(2,3]$ and $\overline{x}\in\!\mathcal{X}^*$. Let $0<\!\varrho<\!\frac{1-\varsigma}{1+\varsigma+L_g}$ for some $\varsigma\in(0,1)$, and $\eta\!:=\!\frac{2+L_HL_{m}^{-1}}{2\eta_0}\!+\!\frac{\sqrt{2L_H(\eta_0L_{m})^{-1}+\eta_0^{-2}(2+L_HL_{m}^{-1})^2}}{2}$ for $\eta_0\!=1-\!\varsigma\!-\!\varrho(1\!+\!\varsigma\!+\!L_g)$. Under Assumptions \ref{ass1}-\ref{ass2}, 
 there is $\widetilde{k}\in\!\mathbb{N}$ such that for all $k\ge\widetilde{k}$ with $x^k\!\in\mathbb{B}(\overline{x},{\overline{\varepsilon}}/{2})$,
 \begin{equation}\label{ykxk-bound}
   \|y^k-x^k\|\le \eta{\rm dist}(x^k,\mathcal{X}^*).
 \end{equation}
\end{proposition}

 \begin{equation}\label{ykxk-relation}
  (1\!-\!\varsigma)\|y^k-x^k\|\le\|y^k-v^k-x^k\|\le(1\!+\!\varsigma)\|y^k-x^k\|.
 \end{equation}

\begin{theorem}\label{Theorem-Lconverge}
 Fix any $q\in(2,3]$ and $\overline{x}\in\!\omega(x^0)$. Let  $0<\varrho<\!\frac{1-\varsigma}{1+\varsigma+L_g}$ for some $\varsigma\in(0,1)$. Suppose that Assumptions \ref{ass1}-\ref{ass2} hold, and that there exist $\widehat{\varepsilon}>0$ and $\kappa>0$ such that for all $x\in\mathbb{B}(\overline{x},\widehat{\varepsilon})$, ${\rm dist}(x,\mathcal{X}^*)\le\kappa[r(x)]^{\gamma}$ with $\gamma\in(\frac{1}{q-1},1]$. 
  Then, with $\widehat{c}\!:=\![(3+\!L_g)\varrho \widehat{L}\!+\!L_H/2\!+\!\widehat{L}]\eta^{q-1}$ where $\eta$ is the same as before, for all $k$ large enough, 
 \begin{equation}\label{aim-rineq}
 r(x^{k+1})\le \kappa\widehat{c}[r(x^k)]^{\gamma(q-1)}\ \ {\rm and}\ \ 
 \|x^{k+1}\!-\!\overline{x}\|
 \le2\eta\kappa(1\!+\!\varsigma)\widehat{c}^{\gamma}\|x^k\!-\!\overline{x}\|^{\gamma(q-1)},
 \end{equation}
 that is, the sequences $\{r(x^k)\}_{k\in\mathbb{N}}$ and $\{x^k\}_{k\in\mathbb{N}}$ converge to $0$ and $\overline{x}\in\mathcal{X}^*$, respectively, with the $Q$-superlinear rate of order $\gamma(q\!-\!1)$.
\end{theorem}

\begin{proof}
 From $\overline{x}\in\omega(x^0)$, there exists $\mathcal{K}\subset\mathbb{N}$ such that $\lim_{\mathcal{K}\ni k\to\infty}x^k=\overline{x}$. Along with the given assumption, ${\rm dist}(x^k,\mathcal{X}^*)\le\kappa[r(x^k)]^{\gamma}$ for all $k\in\mathcal{K}$ large enough. Passing the limit $\mathcal{K}\ni k\to\infty$ and using Proposition \ref{prop2-xk} (ii) leads to $\overline{x}\in\mathcal{X}^*$. 
 For each $k\ge\widetilde{k}$ where $\widetilde{k}$ is the same as in Proposition \ref{prop-dkbound1}, since $x^k\!\in\Gamma$ and $\overline{x}\in\omega(x^0)\!\subset\Gamma$, invoking the nonexpansiveness of ${\rm prox}_g$ and equation \eqref{grad-Lip} results in 
 \begin{equation}\label{rxk-ineq41}
  r(x^k)=|r(x^k)-r(\overline{x})|\le 2\|x^k-\overline{x}\|+\|\nabla\!f(x^k)-\nabla\!f(\overline{x})\|\le (2\!+\!L_g)\|x^k\!-\!\overline{x}\|.
 \end{equation} 
 Let $\varepsilon_1\!:=\min\{\widehat{\varepsilon},\overline{\varepsilon}\}/2$. From \eqref{new-iter} and \eqref{ykxk-relation}, for all $k\ge \widetilde{k}$, it holds that 
 \begin{equation}\label{xkp1-ineq41}
 (1\!-\!\varsigma)\|y^k\!-\!x^k\|\le\|x^{k+1}\!-\!x^k\|\le(1\!+\!\varsigma)\|y^k\!-\!x^k\|.
 \end{equation}
 Along with $\lim_{k\to\infty}(y^k\!-\!x^k)=0$, we get $\lim_{k\to\infty}(x^{k+1}\!-\!x^k)=0$, so $\|x^{k+1}\!-\!x^k\|\le\varepsilon_1$ for all $k\ge\widetilde{k}$ (if necessary by increasing $\widetilde{k}$). Thus, for all $k\ge \widetilde{k}$ with $x^k\in\mathbb{B}(\overline{x},\varepsilon_1)$, $\|x^{k+1}-\overline{x}\|\le2\varepsilon_1\le\min\{\widehat{\varepsilon},\overline{\varepsilon}\}$. Fix any $k\ge \widetilde{k}$ with $x^k\in\mathbb{B}(\overline{x},\varepsilon_1)$. Then,
 \begin{align}\label{ineq0-rk}
  r(x^{k+1})
  &\stackrel{ \eqref{yk-vk}}{=}\|x^{k+1}-{\rm prox}_g(x^{k+1}\!-\!\nabla\!f(x^{k+1}))+v^k-y^k+{\rm prox}_g(y^k\!-\!\nabla\!f_{k,j_k}(y^k))\|\nonumber\\
  &\le \|x^{k+1}-y^k-\nabla\!f(x^{k+1})+\nabla\!f_{k,j_k}(y^k)\|+\|x^{k+1}-y^k\|+\|v^k\|\nonumber\\
  &\le \|\nabla\!f(x^{k+1})-\nabla\vartheta_{k}(y^k)\|+2\|x^{k+1}-y^k\|+\|v_k\|+L_k\|y^k-x^k\|^{q-1}\nonumber\\
  &\stackrel{ \eqref{new-iter}}{\le}3\|v^k\|+\|\nabla\!f(x^{k+1})-\nabla\vartheta_{k}(y^k)\|+L_k\|y^k-x^k\|^{q-1},
  \end{align}
  where the first inequality is using the nonexpansiveness of ${\rm prox}_g$, and the second one is by the expression of $f_{k,j_k}$ in \eqref{subprobk}. By invoking \eqref{grad-Lip} and \eqref{ineq-quad}, we have
  \[
   \|\nabla\!f(x^{k+1})\!-\!\nabla\vartheta_k(y^k)\|
   \le L_g\|v^k\|+({L_H}/{2})\|y^k\!-\!x^k\|^{2}.
  \]
  Combining this inequality with \eqref{ineq0-rk} and using $\|v^k\|\le\varrho L_k\|y^k\!-\!x^k\|^{q-1}$ leads to  
  \begin{align}\label{rk-rate}
   r(x^{k+1})
   &\le (3+\!L_g)\varrho L_k\|y^k\!-\!x^k\|^{q-1}+({L_H}/{2})\|y^k\!-\!x^k\|^{2}+L_k\|y^k\!-\!x^k\|^{q-1}\nonumber\\
   &\le(3+\!L_g)\varrho \widehat{L}\|y^k\!-\!x^k\|^{q-1}+({L_H}/{2})\|y^k\!-\!x^k\|^{q-1}+\widehat{L}\|y^k\!-\!x^k\|^{q-1}\nonumber\\
   &\stackrel{\eqref{ykxk-bound}}{\le}[(3+\!L_g)\varrho \widehat{L}+L_H/2+\widehat{L}]\eta^{q-1}{\rm dist}(x^k,\mathcal{X}^*)^{q-1}\nonumber\\
   &\le[(3+\!L_g)\varrho \widehat{L}+L_H/2+\widehat{L}]\eta^{q-1}\kappa[r(x^k)]^{\gamma(q-1)},
  \end{align}
  where the second inequality is due to $\|y^k\!-\!x^k\|\le 1$ and the boundedness of $L_k$,  and the fourth one is by the given assumption. Fix  any $\sigma\in(0,1)$. Since $\lim_{k\to\infty}r(x^k)=0$ and $\gamma(q\!-\!1)>1$, from the above \eqref{rk-rate}, there exists $\varepsilon_2\in(0,\varepsilon_1)$ such that 
  \begin{equation}\label{rk-recursion}
   r(x^{k+1})\le\sigma r(x^k)\quad{\rm for\ all}\ k\ge \widetilde{k}\ {\rm with}\ x^k\in\mathbb{B}(\overline{x},\varepsilon_2).
  \end{equation}
  Let $\widetilde{\varepsilon}\!:=\min\{\frac{\varepsilon_2}{(1+\varsigma)\eta+1},\big(\frac{(1-\widetilde{\sigma})\varepsilon_2}{2(1+\varsigma)\eta\kappa(2+L_g)^{\gamma}}\big)^{1/\gamma},\frac{\varepsilon_2}{2}\}$ and $\widetilde{\sigma}:=\sigma^{\gamma}\in(0,1)$. 
  Recall that $\overline{x}\in\omega(x^0)$, there exists $\overline{k}_2\ge \widetilde{k}$ such that $x^{\overline{k}_2}\in\mathbb{B}(\overline{x},\widetilde{\varepsilon})$. We argue by induction that 
  \begin{equation}\label{induction}
  x^{k+1}\in\mathbb{B}(\overline{x},\varepsilon_2)\quad{\rm for\ all}\  k\ge \overline{k}_2. 
  \end{equation}
  Indeed, when $k=\overline{k}_2$, the conclusion holds by invoking $x^{\overline{k}_2}\in\mathbb{B}(\overline{x},\widetilde{\varepsilon})$ and noting that
  \begin{align*}
  \|x^{\overline{k}_2+1}-\overline{x}\|
  &\le\|x^{\overline{k}_2+1}-x^{\overline{k}_2}\|+\|x^{\overline{k}_2}-\overline{x}\|\stackrel{\eqref{xkp1-ineq41}}{\le} (1\!+\!\varsigma)\|y^{\overline{k}_2}\!-\!x^{\overline{k}_2}\|+\|x^{\overline{k}_2}\!-\!\overline{x}\|\\
  &\stackrel{\eqref{ykxk-bound}}{\le} (1\!+\!\varsigma)\eta{\rm dist}(x^{\overline{k}_2},\mathcal{X}^*)+\|x^{\overline{k}_2}-\overline{x}\|\le\big((1\!+\!\varsigma)\eta\!+\!1\big)\|x^{\overline{k}_2}\!-\!\overline{x}\|\le\varepsilon_2,  
 \end{align*}
 where the fourth inequality is due to ${\rm dist}(x^{\overline{k}_2},\mathcal{X}^*)\le\|x^{\overline{k}_2}-\overline{x}\|\le\widetilde{\varepsilon}\le1$. To proceed the induction, fix any $k>\overline{k}_2$ and assume that $x^{l+1}\in\mathbb{B}(\overline{x},\varepsilon_2)$ for all $\overline{k}_2\le l\le k\!-\!1$. By applying the given local error bound assumption, we can obtain that  
 \begin{align*}
  \|x^{k+1}\!-\!x^{\overline{k}_2}\|
  &\le{\textstyle\sum_{l=\overline{k}_2}^{k}}\|x^{l+1}-x^l\|
  \stackrel{\eqref{xkp1-ineq41}}{\le} (1\!+\!\varsigma){\textstyle\sum_{l=\overline{k}_2}^{k}}\|y^l\!-\!x^l\|\\
  &\stackrel{\eqref{ykxk-bound}}{\le} \eta(1\!+\!\varsigma){\textstyle\sum_{l=\overline{k}_2}^{k}}{\rm dist}(x^l,\mathcal{X}^*)
  \le \eta\kappa(1\!+\!\varsigma){\textstyle\sum_{l=\overline{k}_2}^{k}}[r(x^l)]^{\gamma}\\ 
  &\stackrel{\eqref{rk-recursion}}{\le} \eta\kappa(1\!+\!\varsigma)[r(x^{\overline{k}_2})]^{\gamma}{\textstyle\sum_{l=\overline{k}_2}^{k}}\widetilde{\sigma}^{l-\overline{k}_2}
  \le\frac{\eta\kappa(1\!+\!\varsigma)}{1-\widetilde{\sigma}}[r(x^{\overline{k}_2})]^{\gamma}\\
  &\stackrel{\eqref{rxk-ineq41}}{\le}\frac{\eta\kappa(1\!+\!\varsigma)}{1-\widetilde{\sigma}}(2\!+\!L_g)^{\gamma}\|x^{\overline{k}_2}-\overline{x}\|^{\gamma}\le\frac{\varepsilon_2}{2}.
  \end{align*}
  This means that $\|x^{k+1}-\overline{x}\|\le \varepsilon_2/2+\|x^{\overline{k}_2}-\overline{x}\|\le\varepsilon_2$. 
   
  Now we are in a position to prove the convergence result. As $\lim_{k\to\infty}r(x^k)=0$, for any $\varepsilon>0$, there exists $\widehat{k}\ge \overline{k}_2$ such that $r(x^k)\le\varepsilon$ for all $k\ge\widehat{k}$. Fix any $i_1>i_2\ge\widehat{k}$. Using the same arguments as those for the above inequality yields that
  \begin{align*}%\label{ineq-cauchy}
   \|x^{i_1}-x^{i_2}\|\le\sum_{l=i_2}^{i_1-1}\|x^{l+1}-x^l\|
    \le\frac{(1+\varsigma)\eta\kappa}{1-\widetilde{\sigma}}[r(x^{i_2})]^{\gamma}
   \le\frac{(1+\varsigma)\eta\kappa}{1-\widetilde{\sigma}}\varepsilon^{\gamma},
  \end{align*}
  which implies that $\{x^k\}_{k\in\mathbb{N}}$ is a Cauchy sequence, so converges to $\overline{x}\in\mathcal{X}^*$. Next we deduce the convergence rate of $\{r(x^k)\}_{k\in\mathbb{N}}$ and $\{x^k\}_{k\in\mathbb{N}}$. Fix any $k>\widehat{k}$. Then $x^k\in\mathbb{B}(\overline{x},\varepsilon_2)$ follows \eqref{induction}, so \eqref{rk-rate} holds for any $k>\widehat{k}$. That is, the first inequality in \eqref{aim-rineq} holds. From the local error bound and the third inequality in \eqref{rk-rate},  
  \begin{equation}\label{dist-rate}
  {\rm dist}(x^{k+1},\mathcal{X}^*)\le\kappa[(3+\!L_g)\varrho \widehat{L}\!+\!{L_H}/{2}\!+\!\widehat{L}]^{\gamma}\eta^{\gamma(q-1)}{\rm dist}(x^k,\mathcal{X}^*)^{\gamma(q-1)}.
 \end{equation}
 Note that $\lim_{k\to\infty}{\rm dist}(x^k,\mathcal{X}^*)=0$ by the local error bound and $\lim_{k\to\infty}r(x^k)=0$. Inequality \eqref{dist-rate} implies that for any $k>\widehat{k}$, ${\rm dist}(x^{k+1},\mathcal{X}^*)\le\frac{1}{2}{\rm dist}(x^k,\mathcal{X}^*)$. Then, 
  \begin{align}\label{ineq-xk}
   \|x^{i_1}\!-\!x^{i_2}\|
   &\le{\textstyle\sum_{k=i_2}^{i_1-1}}\|x^{k+1}\!-\!x^k\|\stackrel{\eqref{xkp1-ineq41}}{\le}
    (1\!+\!\varsigma){\textstyle\sum_{k=i_2}^{i_1-1}}\|y^k\!-\!x^k\|\nonumber\\
   &\le(1\!+\!\varsigma)\eta{\textstyle\sum_{k=i_2}^{i_1-1}}{\rm dist}(x^k,\mathcal{X}^*)\nonumber\\
   &\stackrel{\eqref{ykxk-bound}}{\le}(1\!+\!\varsigma)\eta{\rm dist}(x^{i_2},\mathcal{X}^*){\textstyle\sum_{k=i_2}^{i_1-1}}(1/2)^{k-i_2}\nonumber\\
   &\le 2(1\!+\!\varsigma)\eta{\rm dist}(x^{i_2},\mathcal{X}^*)\quad\ {\rm for\ any}\ i_1>i_2>\widehat{k}.
  \end{align}
 Letting $i_2=k+1$ and passing the limit $i_1\to\infty$ to the both sides of \eqref{ineq-xk} results in
 \begin{align*}
 \|x^{k+1}\!-\!\overline{x}\|
 &\le 2(1\!+\!\varsigma)\eta{\rm dist}(x^{k+1},\mathcal{X}^*)\\
 &\stackrel{\eqref{dist-rate}}{\le}2(1\!+\!\varsigma)\eta\kappa[(3\!+\!L_g)\varrho \widehat{L}\!+\!L_H/2\!+\!\widehat{L}]^{\gamma}\eta^{\gamma(q-1)}{\rm dist}(x^k,\mathcal{X}^*)^{\gamma(q-1)}\\
 &\le2\eta\kappa(1\!+\!\varsigma)[(3\!+\!L_g)\varrho \widehat{L}+L_H/2+\widehat{L}]^{\gamma}\eta^{\gamma(q-1)} \|x^k\!-\!\overline{x}\|^{\gamma(q-1)}.
 \end{align*} 
 This shows that the second inequality in \eqref{aim-rineq} holds. The proof is completed.
 \end{proof}

\end{document}
